% GENERATED FILE. Edit canonical lessons, then run npm run generate:books.
% canonical-source-hash: fnv1a64:ffdcde83b26a0263
% canonical-lessons: SA-C308-samuhah, SA-C308-upakaroti, SA-C308-vibhajati, SA-C308-data, SA-C308-sahayah

\chapter{Lending a Hand}
\label{ch:sa-lending-a-hand}
\hlchaptermodality{308}

\begin{chapteropening}
\textbf{By the end of this chapter:} \emph{I can say who helped and what they shared.}

Complete the last lesson of chapter 308: I can say who helped and what they shared.
\end{chapteropening}

\section[\texorpdfstring{\sk{समूहः} (samūhaḥ)}{समूहः (samūhaḥ)}]{\sk{समूहः} (samūhaḥ) — a group; a crowd}

\label{lesson:SA-C308-samuhah}



\begin{quote}
Before the new one: say the Sanskrit for orphaned, then the Sanskrit for safety.
\end{quote}

\subsection*{You'll want to know: \sk{समूहः}}
\textbf{\sk{समूहः}} — \emph{samūhaḥ} — \textquotedblleft{}a group; a crowd\textquotedblright{}.

\begin{quote}
\sk{एकः} \sk{समूहः} \sk{शिबिरे} \sk{साहाय्यं} \sk{करोति।}
\end{quote}

\begin{quote}
\emph{ekaḥ samūhaḥ śibire sāhāyyaṃ karoti.}
\end{quote}

\begin{quote}
A group helps out in the camp.
\end{quote}

A masculine noun.

\begin{cousinweb}
\textbf{\sk{समूहः}} (\emph{samūhaḥ}) is built from \emph{sam}, \textquotedblleft{}together\textquotedblright{}, and the root √ūh, \textquotedblleft{}to move, push\textquotedblright{}; MW gives \emph{samūha} as a collection, assemblage, multitude (MW).
\end{cousinweb}

\begin{grammarlens}[title={What you've built}]
One more word for telling what happened.
\end{grammarlens}

\subsection*{Your turn}
\begin{itemize}
\raggedright
  \item \emph{Say it:} \emph{samūhaḥ}
  \item \emph{Say it:} \emph{samūhaḥ}, once more
  \item \emph{Say it:} \emph{ekaḥ samūhaḥ śibire sāhāyyaṃ karoti.} — A group helps out in the camp.
  \item \emph{Recall:} say the Sanskrit for orphaned, then the Sanskrit for safety, then say \emph{samūhaḥ} again
\end{itemize}

\begin{tcolorbox}[breakable,colback=teal!4,colframe=teal!35!black,title={Before you move on}]
What does \sk{समूहः} mean? (\textquotedblleft{}A group; a crowd\textquotedblright{}.) Say it once more.
\end{tcolorbox}
\section[\texorpdfstring{\sk{उपकरोति} (upakaroti)}{उपकरोति (upakaroti)}]{\sk{उपकरोति} (upakaroti) — does a kindness; helps out}

\label{lesson:SA-C308-upakaroti}



\begin{quote}
Before the new one: say the Sanskrit for safety, then the Sanskrit for a group.
\end{quote}

\subsection*{You'll want to know: \sk{उपकरोति}}
\textbf{\sk{उपकरोति}} — \emph{upakaroti} — \textquotedblleft{}does a kindness; helps out\textquotedblright{}.

\begin{quote}
\sk{सः} \sk{दरिद्रस्य} \sk{उपकरोति।}
\end{quote}

\begin{quote}
\emph{saḥ daridrasya upakaroti.}
\end{quote}

\begin{quote}
He does a poor man a kindness.
\end{quote}

Past (he …): \sk{सः} \sk{दरिद्रस्य} \sk{उपकृतवान्}, \emph{saḥ daridrasya upakṛtavān}.

\begin{cousinweb}
\textbf{\sk{उपकरोति}} (\emph{upakaroti}) is built from \emph{upa}, \textquotedblleft{}near, towards\textquotedblright{}, and the root √kṛ, \textquotedblleft{}to do\textquotedblright{}; MW gives \emph{upa-kṛ} as to assist, help, favour (MW).
\end{cousinweb}

\begin{grammarlens}[title={What you've built}]
One more word for telling what happened.
\end{grammarlens}

\subsection*{Your turn}
\begin{itemize}
\raggedright
  \item \emph{Say it:} \emph{upakaroti}
  \item \emph{Say it:} \emph{upakaroti}, once more
  \item \emph{Say it:} \emph{saḥ daridrasya upakaroti.} — He does a poor man a kindness.
  \item \emph{Recall:} say the Sanskrit for safety, then the Sanskrit for a group, then say \emph{upakaroti} again
\end{itemize}

\begin{tcolorbox}[breakable,colback=teal!4,colframe=teal!35!black,title={Before you move on}]
What does \sk{उपकरोति} mean? (\textquotedblleft{}Does a kindness; helps out\textquotedblright{}.) Say it once more.
\end{tcolorbox}
\section[\texorpdfstring{\sk{विभजति} (vibhajati)}{विभजति (vibhajati)}]{\sk{विभजति} (vibhajati) — divides; shares out}

\label{lesson:SA-C308-vibhajati}



\begin{quote}
Before the new one: say the Sanskrit for a group, then the Sanskrit for does a kindness.
\end{quote}

\subsection*{You'll want to know: \sk{विभजति}}
\textbf{\sk{विभजति}} — \emph{vibhajati} — \textquotedblleft{}divides; shares out\textquotedblright{}.

\begin{quote}
\sk{ते} \sk{अन्नं} \sk{सर्वेषु} \sk{विभजन्ति।}
\end{quote}

\begin{quote}
\emph{te annaṃ sarveṣu vibhajanti.}
\end{quote}

\begin{quote}
They share the food out among everyone.
\end{quote}

Past (he …): \sk{सः} \sk{अन्नं} \sk{विभक्तवान्}, \emph{saḥ annaṃ vibhaktavān}.

\begin{cousinweb}
\textbf{\sk{विभजति}} (\emph{vibhajati}) is built from \emph{vi}, \textquotedblleft{}apart\textquotedblright{}, and the root √bhaj, \textquotedblleft{}to share\textquotedblright{}, as in \emph{bhajati}; MW gives \emph{vi-bhaj} as to divide, distribute, apportion (MW).
\end{cousinweb}

\begin{grammarlens}[title={What you've built}]
One more word for telling what happened.
\end{grammarlens}

\subsection*{Your turn}
\begin{itemize}
\raggedright
  \item \emph{Say it:} \emph{vibhajati}
  \item \emph{Say it:} \emph{vibhajati}, once more
  \item \emph{Say it:} \emph{te annaṃ sarveṣu vibhajanti.} — They share the food out among everyone.
  \item \emph{Recall:} say the Sanskrit for a group, then the Sanskrit for does a kindness, then say \emph{vibhajati} again
\end{itemize}

\begin{tcolorbox}[breakable,colback=teal!4,colframe=teal!35!black,title={Before you move on}]
What does \sk{विभजति} mean? (\textquotedblleft{}Divides; shares out\textquotedblright{}.) Say it once more.
\end{tcolorbox}
\section[\texorpdfstring{\sk{दाता} (dātā)}{दाता (dātā)}]{\sk{दाता} (dātā) — a giver; a donor}

\label{lesson:SA-C308-data}



\begin{quote}
Before the new one: say the Sanskrit for does a kindness, then the Sanskrit for shares out.
\end{quote}

\subsection*{You'll want to know: \sk{दाता}}
\textbf{\sk{दाता}} — \emph{dātā} — \textquotedblleft{}a giver; a donor\textquotedblright{}.

\begin{quote}
\sk{दाता} \sk{शिबिराय} \sk{धनं} \sk{ददाति।}
\end{quote}

\begin{quote}
\emph{dātā śibirāya dhanaṃ dadāti.}
\end{quote}

\begin{quote}
The donor gives money to the camp.
\end{quote}

A masculine noun.

\begin{cousinweb}
\textbf{\sk{दाता}} (\emph{dātā}) comes from the root √dā, \textquotedblleft{}to give\textquotedblright{}, with the ending \emph{-tṛ} of the one who does it; MW gives \emph{dātṛ} as a giver, donor (MW).
\end{cousinweb}

\begin{grammarlens}[title={What you've built}]
One more word for telling what happened.
\end{grammarlens}

\subsection*{Your turn}
\begin{itemize}
\raggedright
  \item \emph{Say it:} \emph{dātā}
  \item \emph{Say it:} \emph{dātā}, once more
  \item \emph{Say it:} \emph{dātā śibirāya dhanaṃ dadāti.} — The donor gives money to the camp.
  \item \emph{Recall:} say the Sanskrit for does a kindness, then the Sanskrit for shares out, then say \emph{dātā} again
\end{itemize}

\begin{tcolorbox}[breakable,colback=teal!4,colframe=teal!35!black,title={Before you move on}]
What does \sk{दाता} mean? (\textquotedblleft{}A giver; a donor\textquotedblright{}.) Say it once more.
\end{tcolorbox}
\section[\texorpdfstring{\sk{सहायः} (sahāyaḥ)}{सहायः (sahāyaḥ)}]{\sk{सहायः} (sahāyaḥ) — a helper; a companion}

\label{lesson:SA-C308-sahayah}



\begin{quote}
Before the new one: say the Sanskrit for shares out, then the Sanskrit for a donor.
\end{quote}

\subsection*{You'll want to know: \sk{सहायः}}
\textbf{\sk{सहायः}} — \emph{sahāyaḥ} — \textquotedblleft{}a helper; a companion\textquotedblright{}.

\begin{quote}
\sk{अरुणः} \sk{वैद्यस्य} \sk{सहायः} \sk{अस्ति।}
\end{quote}

\begin{quote}
\emph{aruṇaḥ vaidyasya sahāyaḥ asti.}
\end{quote}

\begin{quote}
Arun is the doctor's helper.
\end{quote}

A masculine noun.

\begin{cousinweb}
\textbf{\sk{सहायः}} (\emph{sahāyaḥ}) MW thinks it is probably built from \emph{saha}, \textquotedblleft{}with\textquotedblright{}, and \emph{aya}, \textquotedblleft{}going\textquotedblright{}, one who goes along with another, though some take it as a Prakrit form, so its origin is not settled (MW).
\end{cousinweb}

\begin{grammarlens}[title={What you've built}]
One more word for telling what happened.
\end{grammarlens}

\subsection*{Your turn}
\begin{itemize}
\raggedright
  \item \emph{Say it:} \emph{sahāyaḥ}
  \item \emph{Say it:} \emph{sahāyaḥ}, once more
  \item \emph{Say it:} \emph{aruṇaḥ vaidyasya sahāyaḥ asti.} — Arun is the doctor's helper.
  \item \emph{Recall:} say the Sanskrit for shares out, then the Sanskrit for a donor, then say \emph{sahāyaḥ} again
\end{itemize}

\begin{tcolorbox}[breakable,colback=teal!4,colframe=teal!35!black,title={Before you move on}]
What does \sk{सहायः} mean? (\textquotedblleft{}A helper; a companion\textquotedblright{}.) Say it once more.
\end{tcolorbox}
